% ID: SEL-D0069-OB-E01
% ESTADO: validada. S069.
\textbf{a)} El recinto se sitúa entre la cota superior $y=4$ y las inferiores $y=x$ e $y=2-x$, con $x\ge0$. En $x=0$ se obtienen $(0,2)$ y $(0,4)$. Las cotas inferiores se cruzan en $(1,1)$ y $y=x$ alcanza la cota superior en $(4,4)$. Los cuatro puntos delimitan el recinto cerrado.
\begin{center}
\setlength{\unitlength}{0.9mm}
\begin{picture}(89,63)(-17,-8)
\qbezier(0,0)(34,0)(68,0)
\qbezier(0,0)(0,25)(0,50)
\qbezier(0,22)(6,16.5)(12,11)
\qbezier(12,11)(30,27.5)(48,44)
\qbezier(48,44)(24,44)(0,44)
\put(69,0){$x$}\put(2,50){$y$}
\put(-15,21){\scriptsize $(0,2)$}\put(-15,44){\scriptsize $(0,4)$}
\put(9,5){\scriptsize $(1,1)$}\put(49,44){\scriptsize $(4,4)$}
\end{picture}
\end{center}
\textbf{b)} Por la primera restricción, $F\ge2$. Hay igualdad en todo el segmento entre $(0,2)$ y $(1,1)$. Como $x\le y\le4$, también se cumple $F\le8$, con igualdad sólo en $(4,4)$. El mínimo es $\boxed{2}$ y el máximo es $\boxed{8}$.

\textbf{c)} En el punto propuesto, $x+y=5/3<2$. Por tanto, queda fuera.
